\documentclass[a4paper,11pt,reqno]{amsart}

% =========================================================================
%  Instrumentación Electrónica — PostLab de la sesión 1 (amplificadores operacionales)
%  Mismo formato que el PreLab de la sesión 1. Los datos medidos se rellenan en los
%  huecos marcados con \dato y \hueco. Compilar con: latexmk -pdf main.tex
% =========================================================================
% =========================================================================
%  Instrumentación Electrónica — PostLab de la sesión 1 (amplificadores operacionales)
%  Estilo heredado de los apuntes de Electromagnetismo I (Temas 0 y 1):
%  amsart, Palatino, versalitas, enunciados numerados por sección, paleta de
%  azules y grises sobrios. Ecuaciones numeradas a la derecha (opción reqno).
% =========================================================================

% --- IDIOMA Y CODIFICACIÓN ------------------------------------------------
\usepackage[utf8]{inputenc}
\usepackage[T1]{fontenc}
\usepackage[spanish,es-tabla,es-nodecimaldot]{babel}

% --- TIPOGRAFÍA (Palatino) --------------------------------------------------
\usepackage{mathpazo}
\linespread{1.05}
\emergencystretch=2em
\usepackage[final]{microtype}

% --- MATEMÁTICAS -------------------------------------------------------------
\usepackage{amsmath, amssymb, mathtools}
\usepackage{bm}

% --- UNIDADES (coma decimal, como en el resto de los apuntes) ---------------
\usepackage{siunitx}
\sisetup{output-decimal-marker={,}, group-separator={\,}, group-minimum-digits=5,
         exponent-product=\cdot, per-mode=symbol, range-phrase={--},
         list-final-separator={ y }, list-pair-separator={ y }}

% --- PÁGINA, COLORES Y GRÁFICOS --------------------------------------------
\usepackage[a4paper,margin=2.6cm]{geometry}
\usepackage{graphicx}
\usepackage{xcolor}
% Paleta base (idéntica a la de los apuntes de Electromagnetismo I)
\definecolor{azulprof}{HTML}{1F3A5F}    % azul oscuro
\definecolor{azulacero}{HTML}{3B6EA5}   % azul medio
\definecolor{azulclaro}{HTML}{8FB3D9}   % azul claro
\definecolor{azulfondo}{HTML}{F3F7FB}   % fondo tenue de definiciones
\definecolor{grispizarra}{HTML}{6B7280} % gris
\definecolor{grisclaro}{HTML}{D4D8DE}   % gris claro
\definecolor{acento}{HTML}{B8800A}      % acento: ámbar oscuro
\definecolor{rojoteo}{HTML}{B8800A}     % señales de salida, resultados
\definecolor{verdeteo}{HTML}{2F6F68}    % corrientes y sentidos de referencia
\definecolor{opfondo}{HTML}{E4EEF8}     % relleno del triángulo del AO

\usepackage{tikz}
\usetikzlibrary{babel, arrows.meta, calc, positioning, decorations.markings,
                patterns, intersections}
\usepackage[european, straightvoltages, RPvoltages]{circuitikz}
\usepackage{pgfplots}
\pgfplotsset{compat=1.17}
\usepgfplotslibrary{fillbetween}
\usepackage{booktabs}
\usepackage{needspace}
\usepackage{array}
\usepackage{multirow}
\usepackage[inline]{enumitem}
\setlist[itemize]{leftmargin=1.6em, itemsep=2pt, topsep=4pt}
\setlist[enumerate]{leftmargin=1.8em, itemsep=2pt, topsep=4pt}

% --- ENLACES Y REFERENCIAS CRUZADAS ----------------------------------------
\usepackage[colorlinks=true, linkcolor=azulacero, citecolor=azulacero,
            urlcolor=azulacero, plainpages=false, pdfusetitle=false,
            pdftitle={PostLab. Sesión 1: amplificadores operacionales},
            pdfauthor={Daniel Martín Guerrero}]{hyperref}
\usepackage[spanish,nameinlink]{cleveref}

% --- CABECERAS ---------------------------------------------------------------
\pagestyle{headings}
\makeatletter
\def\ps@headings{%
  \let\@mkboth\markboth
  \def\@evenhead{\small\scshape\normalfont\thepage\hfil Instrumentación Electrónica}%
  \def\@oddhead{\small\scshape\normalfont PostLab. Sesión 1: amplificadores operacionales\hfil\thepage}%
  \let\@evenfoot\@empty \let\@oddfoot\@empty}
\makeatother

% Subsecciones como títulos en bloque (no «run-in»), con aire antes y después
\usepackage{titlesec}
\titleformat{\subsection}[block]{\normalfont\bfseries\color{azulprof}}{\thesubsection.}{0.6em}{}
\titlespacing*{\subsection}{0pt}{3ex plus 1ex minus .5ex}{1.6ex plus .4ex}
\setcounter{tocdepth}{2}
\makeatletter
\def\@@and{y}
\renewcommand{\keywordsname}{Palabras clave}
% Índice: número de sección en una caja de anchura fija (evita el solape de amsart)
\renewcommand{\tocsection}[3]{\indentlabel{\@ifnotempty{#2}{\makebox[1.6em][l]{#2.}}}#3}
\renewcommand{\tocsubsection}[3]{\indentlabel{\@ifnotempty{#2}{\makebox[2.4em][l]{#2.}}}#3}
\def\l@section{\@tocline{1}{2pt}{0pt}{1.6em}{}}
\def\l@subsection{\@tocline{2}{0pt}{1.6em}{2.4em}{}}
\def\l@part{\@tocline{-1}{8pt plus 2pt}{0pt}{0pt}{\bfseries\color{azulprof}}}
\makeatother

% =========================================================================
%               NOTACIÓN Y MACROS DE USO FRECUENTE
% =========================================================================
\newcommand{\dif}{\mathrm{d}}
\newcommand{\eul}{\mathrm{e}}
\newcommand{\jj}{\mathrm{j}}                         % unidad imaginaria (j)
\newcommand{\R}{\mathbb{R}}
\newcommand{\C}{\mathbb{C}}
\newcommand{\abs}[1]{\lvert #1\rvert}
\newcommand{\defeq}{\mathrel{\overset{\mathrm{def}}{=}}}
\newcommand{\Lt}{\mathcal{L}}                        % transformada de Laplace
\DeclareMathOperator{\sen}{sen}
\DeclareMathOperator{\sgn}{sgn}
\DeclareMathOperator{\arctg}{arctg}
\DeclareMathOperator{\Real}{Re}
\renewcommand{\sin}{\sen}                            % «sen» en todo el texto
% Magnitudes de circuito
\newcommand{\vin}{v_{\mathrm{in}}}
\newcommand{\vout}{v_{\mathrm{out}}}
\newcommand{\Vin}{V_{\mathrm{in}}}
\newcommand{\Vout}{V_{\mathrm{out}}}
\newcommand{\Vsat}{V_{\mathrm{sat}}}
\newcommand{\Vcc}{V_{\mathrm{CC}}}
\newcommand{\Vee}{V_{\mathrm{EE}}}
\newcommand{\Rin}{R_{\mathrm{in}}}
\newcommand{\Rout}{R_{\mathrm{out}}}
\newcommand{\Zin}{Z_{\mathrm{in}}}
\newcommand{\Zout}{Z_{\mathrm{out}}}
\newcommand{\Av}{A_{V}}
\newcommand{\vp}{v_{+}}
\newcommand{\vm}{v_{-}}
\newcommand{\dB}{\,\mathrm{dB}}
\newcommand{\pend}[1]{\textcolor{grispizarra}{\colorbox{grisclaro}{\strut\,#1\,}}}

% =========================================================================
%        ENTORNOS FORMALES (enunciados numerados, estilo artículo)
% =========================================================================
\usepackage{amsthm}
\usepackage[most]{tcolorbox}

\newtheoremstyle{cienplano}{10pt}{10pt}{\itshape}{}%
  {\bfseries\color{azulprof}}{.}{0.6em}%
  {\thmname{#1}\thmnumber{ #2}\thmnote{\normalfont\upshape\ (#3)}}
\newtheoremstyle{ciendef}{10pt}{10pt}{\normalfont}{}%
  {\bfseries\color{azulprof}}{.}{0.6em}%
  {\thmname{#1}\thmnumber{ #2}\thmnote{\normalfont\upshape\ (#3)}}
\newtheoremstyle{cienobs}{8pt}{8pt}{\normalfont}{}%
  {\itshape\color{azulprof}}{.}{0.6em}%
  {\thmname{#1}\thmnumber{ #2}\thmnote{\normalfont\upshape\ (#3)}}

\theoremstyle{cienplano}
\newtheorem{proposicion}{Proposición}
\newtheorem{teorema}[proposicion]{Teorema}
\newtheorem{corolario}[proposicion]{Corolario}
\theoremstyle{ciendef}
\newtheorem{ejemplo}[proposicion]{Ejemplo}
\newtheorem{ejercicio}[proposicion]{PreLab}
\newtheorem{hipotesis}[proposicion]{Hipótesis}
% Definiciones: único entorno enmarcado (caja fina, fondo tenue, título en el borde)
\newcounter{definicion}
\renewcommand{\thedefinicion}{\theproposicion}
\newtcolorbox{definicionbox}[1]{%
  enhanced, breakable, sharp corners,
  colback=azulfondo, colframe=azulclaro, boxrule=0.4pt,
  left=9pt, right=9pt, top=5pt, bottom=5pt,
  before skip=12pt, after skip=12pt,
  attach boxed title to top left={yshift*=-\tcboxedtitleheight/2, xshift=8pt},
  boxed title style={colback=white, colframe=white, sharp corners, boxrule=0pt,
                     left=3pt, right=3pt, top=1pt, bottom=1pt},
  fonttitle=\bfseries\color{azulprof}, coltitle=azulprof,
  title={Definición~\thedefinicion\ifx\relax#1\relax\else\ (#1)\fi}}
\newenvironment{definicion}[1][]
  {\par\stepcounter{proposicion}\refstepcounter{definicion}\begin{definicionbox}{#1}}
  {\end{definicionbox}}

% Enunciado de cada recuadro del PreLab: misma caja, con filete ámbar (los
% recuadros amarillos del guion), para distinguir lo pedido de lo resuelto.
\newcounter{prelab}
\newtcolorbox{prelabbox}[1]{%
  enhanced, sharp corners,
  colback=white, colframe=acento!70, boxrule=0.5pt, leftrule=2.2pt,
  left=9pt, right=9pt, top=5pt, bottom=5pt,
  before skip=12pt, after skip=10pt,
  attach boxed title to top left={yshift*=-\tcboxedtitleheight/2, xshift=8pt},
  boxed title style={colback=white, colframe=white, sharp corners, boxrule=0pt,
                     left=3pt, right=3pt, top=1pt, bottom=1pt},
  fonttitle=\bfseries\color{acento!85!black}, coltitle=acento!85!black,
  title={PreLab~\theprelab\ifx\relax#1\relax\else\ (#1)\fi}}
\newenvironment{prelab}[1][]
  {\par\refstepcounter{prelab}\begin{prelabbox}{#1}}
  {\end{prelabbox}}

\theoremstyle{cienobs}
\newtheorem{observacion}[proposicion]{Observación}

% Resolución: misma maquinaria que «proof», con otro rótulo
\renewcommand{\qedsymbol}{\ensuremath{\blacksquare}}
\newenvironment{resolucion}[1][Resolución]
  {\begin{proof}[\itshape\color{azulprof}#1]}{\end{proof}}

% Resultado final enmarcado (como \boxed, pero con el filete de la paleta)
\newcommand{\resultado}[1]{\fcolorbox{azulacero}{white}{\ensuremath{\displaystyle #1}}}

\crefname{proposicion}{proposición}{proposiciones}
\Crefname{proposicion}{Proposición}{Proposiciones}
\crefname{teorema}{teorema}{teoremas}
\Crefname{teorema}{Teorema}{Teoremas}
\crefname{corolario}{corolario}{corolarios}
\Crefname{corolario}{Corolario}{Corolarios}
\crefname{definicion}{definición}{definiciones}
\Crefname{definicion}{Definición}{Definiciones}
\crefname{ejemplo}{ejemplo}{ejemplos}
\Crefname{ejemplo}{Ejemplo}{Ejemplos}
\crefname{hipotesis}{hipótesis}{hipótesis}
\Crefname{hipotesis}{Hipótesis}{Hipótesis}
\crefname{prelab}{PreLab}{PreLabs}
\Crefname{prelab}{PreLab}{PreLabs}
\crefname{observacion}{observación}{observaciones}
\Crefname{observacion}{Observación}{Observaciones}
\crefname{section}{sección}{secciones}
\Crefname{section}{Sección}{Secciones}
\crefname{subsection}{sección}{secciones}
\Crefname{subsection}{Sección}{Secciones}
\crefname{table}{tabla}{tablas}
\Crefname{table}{Tabla}{Tablas}
\crefname{figure}{figura}{figuras}
\Crefname{figure}{Figura}{Figuras}
\crefname{equation}{ecuación}{ecuaciones}
\Crefname{equation}{Ecuación}{Ecuaciones}
\crefformat{equation}{ecuación~(#2#1#3)}
\Crefformat{equation}{Ecuación~(#2#1#3)}
\crefrangeformat{equation}{ecuaciones~(#3#1#4)--(#5#2#6)}
\crefmultiformat{equation}{ecuaciones~(#2#1#3)}{ y~(#2#1#3)}{, (#2#1#3)}{ y~(#2#1#3)}

\usepackage[font=small, labelfont={sc,color=azulprof}, labelsep=period,
            margin=0.6cm, justification=justified]{caption}
\usepackage{subcaption}

% Referencia a un panel de figura: «figura 5\,(b)», enlazada en bloque
\newcommand{\figpanel}[2]{\hyperref[#1]{figura~\ref*{#1}\,\textup{(#2)}}}
\newcommand{\Figpanel}[2]{\hyperref[#1]{Figura~\ref*{#1}\,\textup{(#2)}}}

% =========================================================================
%        ESTILOS GRÁFICOS COMUNES (TikZ / circuitikz / pgfplots)
% =========================================================================
% =========================================================================
%  Estilos gráficos comunes: TikZ, circuitikz y pgfplots
% =========================================================================

% --- Estilos TikZ generales ------------------------------------------------
\tikzset{
  >=Stealth,
  eje/.style={->, line width=0.6pt, grispizarra},
  fino/.style={line width=0.4pt, grispizarra},
  punteado/.style={densely dotted, line width=0.5pt, grispizarra},
  etiq/.style={font=\small, text=black!85},
  nota/.style={font=\scriptsize, text=grispizarra, align=center},
  % corriente dibujada junto a un hilo (flecha verde con rótulo)
  corr/.style={-{Stealth[length=4.5pt,width=4pt]}, line width=0.8pt, verdeteo},
  % nudo con tensión anotada
  tens/.style={font=\small, text=azulprof},
}

% --- circuitikz: trazo y colores sobrios -----------------------------------
\ctikzset{
  bipoles/thickness=1,
  bipoles/length=1.15cm,
  tripoles/op amp/width=1.5,
  tripoles/op amp/height=1.35,
  tripoles/op amp/font=\small,
  current/distance=0.45,
  label/align=smart,
}
\tikzset{
  circuito/.style={line width=0.55pt, color=black!85, font=\small},
  ao/.style={fill=opfondo},
}

% --- Estilo base de ejes pgfplots ------------------------------------------
\pgfplotsset{
  tema/.style={
    tick label style={font=\scriptsize}, label style={font=\small},
    axis lines=left, axis line style={grispizarra},
    grid=major, grid style={grisclaro, line width=0.4pt},
    every tick/.style={grispizarra},
    legend style={font=\scriptsize, draw=none, fill=white, fill opacity=0.85,
                  text opacity=1},
    legend cell align=left,
    /pgf/number format/use comma,
  },
  % ejes centrados (formas de onda y características de transferencia)
  temacentro/.style={
    tick label style={font=\scriptsize}, label style={font=\small},
    axis lines=middle, axis line style={grispizarra, -{Stealth[length=4pt]}},
    every tick/.style={grispizarra},
    legend style={font=\scriptsize, draw=none, fill=white, fill opacity=0.85,
                  text opacity=1},
    legend cell align=left,
    /pgf/number format/use comma,
  },
  senal entrada/.style={azulacero, line width=1.0pt},
  senal salida/.style={rojoteo, line width=1.1pt},
  asintota/.style={grispizarra, densely dashed, line width=0.6pt},
}


% =========================================================================
%  Entornos y macros propios del PostLab
% =========================================================================
% Pregunta del guion (recuadros azules): caja con filete azul
\newcounter{pregunta}
\newtcolorbox{preguntabox}[1]{%
  enhanced, sharp corners,
  colback=azulfondo, colframe=azulclaro, boxrule=0.5pt, leftrule=2.2pt,
  left=9pt, right=9pt, top=5pt, bottom=5pt,
  before skip=12pt, after skip=10pt,
  attach boxed title to top left={yshift*=-\tcboxedtitleheight/2, xshift=8pt},
  boxed title style={colback=white, colframe=white, sharp corners, boxrule=0pt,
                     left=3pt, right=3pt, top=1pt, bottom=1pt},
  fonttitle=\bfseries\color{azulprof}, coltitle=azulprof,
  title={Cuestión~\thepregunta\ifx\relax#1\relax\else\ (#1)\fi}}
\newenvironment{pregunta}[1][]
  {\par\refstepcounter{pregunta}\begin{preguntabox}{#1}}
  {\end{preguntabox}}
\newenvironment{respuesta}[1][Respuesta]
  {\begin{proof}[\itshape\color{azulprof}#1]}{\end{proof}}
% Hueco para un oscilograma o gráfica con datos medidos
\newcommand{\hueco}[2]{%
  \tikz\node[draw=azulclaro, densely dashed, line width=0.6pt, fill=azulfondo,
             minimum width=#1, minimum height=3.6cm, align=center,
             font=\small, text=grispizarra, inner sep=4pt]{#2};}
% Dato pendiente de medir
\newcommand{\dato}{\textcolor{grisclaro}{\rule[-0.35ex]{1.3cm}{2.2ex}}}


\newcommand{\Universidad}{Universidad de Salamanca}
\newcommand{\Facultad}{Facultad de Ciencias}
\newcommand{\Titulacion}{Grado en Física}
\newcommand{\Asignatura}{Instrumentación Electrónica}
\newcommand{\NumTema}{Laboratorio de circuitos, sesión 1}
\newcommand{\TituloTema}{PostLab:\\amplificadores operacionales}
\newcommand{\tipotrabajo}{PostLab. Informe de la sesión}
\newcommand{\miNombre}{Daniel Martín Guerrero}
\newcommand{\miRotuloProfesor}{Profesores}
\newcommand{\miProfesor}{Yahya Moubarak Meziani\\María Susana Pérez Santos\\Jesús Enrique Velázquez Pérez}
\newcommand{\miDetalleProfesor}{Área de Electrónica}
\newcommand{\SubtituloTema}{Amplificadores inversor y no inversor,\\integrador y derivador inversores}
\newcommand{\miUbicacion}{Salamanca, curso 2026/2027}

\title{PostLab de la sesión 1: amplificadores operacionales}
\author{Daniel Martín Guerrero}
\address{Grado en Física, Facultad de Ciencias, Universidad de Salamanca,
         Salamanca, España}
\date{}

\begin{document}
\pagenumbering{alph}
%%%%%%%%%%%%%%%%%%%%%%%%%%%%%%%%%%%%%%%%%%%%%%%%%%%%%%%%%%%%%%%%%%%%%%%%
% Portada formal (estilo tesis): logotipos, identificación institucional,
% título y datos de autoría. Banda vertical a la derecha, como en el
% «Cálculo integral multivariable» de L. Navas. Requiere en main.tex las
% macros \Universidad, \Facultad, \Titulacion, \Asignatura, \NumTema,
% \TituloTema, \SubtituloTema, \tipotrabajo, \miNombre, \miProfesor,
% \miUbicacion, \miRotuloProfesor, \miDetalleProfesor.
%%%%%%%%%%%%%%%%%%%%%%%%%%%%%%%%%%%%%%%%%%%%%%%%%%%%%%%%%%%%%%%%%%%%%%%%

\begin{titlepage}
\thispagestyle{empty}
\pagecolor{white}

\begin{tikzpicture}[remember picture, overlay]
  \coordinate (O) at (current page.south west);

  % Velo muy tenue en la mitad superior (cabecera institucional)
  \shade[top color=azulfondo, bottom color=white]
    ($(O)+(0,18.5cm)$) rectangle ($(O)+(21cm,29.7cm)$);

  % Banda vertical azul, a la derecha del bloque de texto
  \fill[azulacero] ($(O)+(19.0cm,2.6cm)$) rectangle ($(O)+(19.14cm,27.4cm)$);

  % Logotipos
  \node[anchor=north west, inner sep=0] at ($(O)+(2.4cm,27.7cm)$)
    {\includegraphics[height=1.55cm]{logo_usal.png}};
  \node[anchor=north east, inner sep=0] at ($(O)+(17.4cm,27.7cm)$)
    {\includegraphics[height=1.55cm]{logo_facultad.png}};
  \draw[grispizarra!60, line width=0.4pt] ($(O)+(2.4cm,25.7cm)$) -- ($(O)+(17.4cm,25.7cm)$);

  % Identificación institucional
  \node[anchor=north, align=center, inner sep=0, text width=15cm] at ($(O)+(9.9cm,25.2cm)$)
    {{\color{azulprof}\scshape\Large\textls[120]{\Universidad}}\\[5pt]
     {\color{grispizarra}\scshape\large\textls[80]{\Facultad}}\\[3pt]
     {\color{grispizarra}\scshape\large\textls[80]{\Titulacion}}};

  % Asignatura y título
  \node[anchor=north, align=center, inner sep=0, text width=15cm] at ($(O)+(9.9cm,19.2cm)$)
    {{\color{azulacero}\scshape\Large\textls[160]{\Asignatura}}\\[4pt]
     {\color{grispizarra}\scshape\large\textls[100]{\NumTema}}\\[22pt]
     {\color{azulprof}\fontsize{30}{36}\selectfont\bfseries\TituloTema}\\[14pt]
     {\color{azulacero}\rule{3.2cm}{0.8pt}}\\[14pt]
     {\color{black!80}\itshape\large\SubtituloTema}};

  % Epígrafe
  \node[anchor=north, align=center, inner sep=0, text width=15cm] at ($(O)+(9.9cm,11.0cm)$)
    {\color{grispizarra}\small
     $\displaystyle \vout(t)=-\frac{1}{RC}\int_{0}^{t}\vin(\tau)\,\dif\tau\qquad
      \vout(t)=-RC\,\frac{\dif\vin}{\dif t}\qquad
      A_V=\frac{\partial\vout}{\partial\vin}$};

  % Autoría y dirección
  \draw[grispizarra!60, line width=0.4pt] ($(O)+(2.4cm,8.6cm)$) -- ($(O)+(17.4cm,8.6cm)$);
  \node[anchor=north west, inner sep=0, text width=7cm, align=left] at ($(O)+(2.4cm,8.2cm)$)
    {{\scshape\small\color{azulacero}\textls[80]{Autor}}\\[4pt]
     {\large\color{black!85}\miNombre}\\[2pt]
     {\small\color{grispizarra}Grado en Física}};
  \node[anchor=north west, inner sep=0, text width=7cm, align=left] at ($(O)+(10.4cm,8.2cm)$)
    {{\scshape\small\color{azulacero}\textls[80]{\miRotuloProfesor}}\\[4pt]
     {\color{black!85}\miProfesor}\\[3pt]
     {\small\color{grispizarra}\miDetalleProfesor}};

  % Tipo de documento, lugar y curso
  \node[anchor=south, align=center, inner sep=0, text width=15cm] at ($(O)+(9.9cm,2.9cm)$)
    {{\color{azulprof}\scshape\large\textls[100]{\tipotrabajo}}\\[4pt]
     {\color{grispizarra}\scshape\small\textls[80]{\miUbicacion}}};
\end{tikzpicture}

\end{titlepage}

\clearpage
\pagenumbering{arabic}
\setcounter{page}{1}

\maketitle

% =========================================================================
\section{Introducción}\label{sec:introduccion}
% =========================================================================
En esta sesión montamos con un $\mu$A741, alimentado a $\pm10\,\mathrm{V}$, los circuitos diseñados en el PreLab:
un amplificador inversor de ganancia $-10$, uno no inversor de ganancia $+11$ y, después, un integrador y un
derivador inversores. Con el osciloscopio medimos las formas de onda de entrada y salida y las
características de transferencia en modo X--Y, y las comparamos con lo que predice el modelo ideal del
amplificador operacional. Este informe recoge esas medidas y responde, en orden, a las cuestiones del guion.

\begin{center}\small
\renewcommand{\arraystretch}{1.3}
\begin{tabular}{@{}lll@{}}
\toprule
Integrante & Apellidos, nombre & DNI/Pasaporte\\
\midrule
1 & \pend{\phantom{Apellido Apellido, Nombre}} & \pend{\phantom{00000000X}}\\
2 & \pend{\phantom{Apellido Apellido, Nombre}} & \pend{\phantom{00000000X}}\\
\bottomrule
\end{tabular}
\end{center}

\textbf{Componentes.} Como pide el guion, medimos con el multímetro el valor real de cada componente pasivo;
los cálculos de este informe usan los valores medidos.
\begin{center}\small
\renewcommand{\arraystretch}{1.3}
\begin{tabular}{@{}lcc@{}}
\toprule
Componente & Valor nominal & Valor medido\\
\midrule
$R_1$ (circuitos 1 y 2) & $1\,\mathrm{k}\Omega$ & \dato\\
$R_2$ (circuitos 1 y 2) & $10\,\mathrm{k}\Omega$ & \dato\\
$R_L$ & $10\,\mathrm{k}\Omega$ & \dato\\
$R$ (circuito 3) & $1\,\mathrm{k}\Omega$ & \dato\\
$R_F$ (circuito 3, set 1) & $100\,\mathrm{k}\Omega$ & \dato\\
$C$ (circuito 3) & $1{,}59\,\mathrm{nF}$ & \dato\\
Tensiones de alimentación $\pm\Vcc$ & $\pm10\,\mathrm{V}$ & \dato\\
\bottomrule
\end{tabular}
\end{center}

% =========================================================================
\section{Circuito 1: amplificador inversor}\label{sec:inversor}
% =========================================================================
Montamos el inversor de la \cref{fig:inversor}, con $R_2/R_1=10$, y lo excitamos con una sinusoide de
$1\,\mathrm{kHz}$ (canal X: $\vin$; canal Y: $\vout$).

% Figura: amplificador inversor (circuito 1) con tensiones y corrientes del análisis.
\begin{figure}[ht]
\centering
\begin{circuitikz}[circuito]
  \draw (0,0) node[op amp, ao] (A) {};
  % nudo inversor y red de realimentación
  \draw (A.-) -- ++(-0.6,0) coordinate (N);
  \draw (N) to[R, l_=$R_1$, i<^=$i_1$] ++(-2.4,0) coordinate (E);
  \draw (N) |- ++(0,1.5) coordinate (T);
  \draw (A.out) -- ++(0.7,0) coordinate (O);
  \draw (T) to[R, l=$R_2$, i>^=$i_2$] (T -| O) -- (O);
  % generador
  \draw (E) to[short, -o] ++(-0.3,0) coordinate (Ein) node[above, tens] {$\vin$};
  \draw (Ein) -- ++(-0.9,0) coordinate (G) to[sV, l_=$v_g$] ++(0,-2.0) node[ground] {};
  % entrada no inversora a tierra
  \draw (A.+) -- ++(-0.6,0) coordinate (P) -- ++(0,-0.45) node[ground] {};
  % carga y salida
  \draw (O) to[short, *-o] ++(1.0,0) coordinate (S) node[right, tens] {$\vout$};
  \draw (O) ++(0.5,0) to[R, l=$R_L$] ++(0,-2.0) node[ground] {};
  \fill (N) circle (1.6pt);
  % anotaciones del análisis
  \node[tens, anchor=east, font=\scriptsize] at ($(N)+(-0.08,0.75)$) {$\vm=0$};
  \draw[azulprof, line width=0.3pt] ($(N)+(-0.08,0.72)$) -- ($(N)+(-0.02,0.06)$);
  \node[tens, anchor=east, font=\scriptsize] at ($(P)+(-0.1,-0.25)$) {$\vp=0$};
\end{circuitikz}
\caption{Amplificador inversor (figura~2 del guion), con $R_1=1\,\mathrm{k}\Omega$,
$R_L=10\,\mathrm{k}\Omega$ y la alimentación omitida. Se indican las magnitudes del análisis: el nudo
inversor es una tierra virtual ($\vm=\vp=0$) y no absorbe corriente, de modo que $i_1=i_2$.}
\label{fig:inversor}
\end{figure}


\begin{pregunta}[Formas de onda y ganancia]
Dibujar en cada caso las formas de onda de $\Vin$ y $\Vout$, detallando las escalas de cada canal y de la base
de tiempos, y comparar la relación de amplitudes (o de tensiones pico a pico) con la $A_V$ buscada.
\end{pregunta}

\begin{respuesta}
Con el modelo ideal, la ganancia esperada, con los valores \emph{medidos} de las resistencias, es
\begin{equation}\label{eq:i-Avteo}
  A_V^{\mathrm{teo}}=-\frac{R_2}{R_1} ,
\end{equation}
y la experimental se obtiene del cociente de tensiones pico a pico, con el signo de la contrafase observada:
\begin{equation}\label{eq:i-Avexp}
  A_V^{\mathrm{exp}}=-\frac{V_{\mathrm{out},pp}}{V_{\mathrm{in},pp}},\qquad
  \varepsilon=\frac{\abs{A_V^{\mathrm{exp}}}-\abs{A_V^{\mathrm{teo}}}}{\abs{A_V^{\mathrm{teo}}}} .
\end{equation}
\begin{center}\small
\renewcommand{\arraystretch}{1.3}
\begin{tabular}{@{}cccccc@{}}
\toprule
$V_p$ nominal & Escala X, Y (V/div) & Base de tiempos & $V_{\mathrm{in},pp}$ & $V_{\mathrm{out},pp}$ & $A_V^{\mathrm{exp}}$\\
\midrule
$0{,}1\,\mathrm{V}$ & \dato & \dato & \dato & \dato & \dato\\
$0{,}2\,\mathrm{V}$ & \dato & \dato & \dato & \dato & \dato\\
$0{,}5\,\mathrm{V}$ & \dato & \dato & \dato & \dato & \dato\\
$1{,}0\,\mathrm{V}$ & \dato & \dato & \dato & \dato & \dato\\
$1{,}5\,\mathrm{V}$ & \dato & \dato & \dato & \dato & \dato\\
\bottomrule
\end{tabular}
\end{center}
\begin{center}
\hueco{0.47\textwidth}{Oscilograma $\vin$, $\vout$\\($V_p=0{,}5\,\mathrm{V}$)}\hfill
\hueco{0.47\textwidth}{Oscilograma $\vin$, $\vout$\\($V_p=1{,}5\,\mathrm{V}$)}
\end{center}
Con $A_V^{\mathrm{teo}}=$ \dato, la desviación \eqref{eq:i-Avexp} es \dato. Mientras la salida no satura, la
salida es una sinusoide en contrafase con la entrada y el cociente pico a pico reproduce \eqref{eq:i-Avteo};
las diferencias se explican por la tolerancia de las resistencias y la resolución de la lectura en pantalla.
La ganancia finita del AO no es apreciable: con $A_0\simeq2\cdot10^{5}$, la ganancia real
$\abs{A_V}=(R_2/R_1)/[1+(1+R_2/R_1)/A_0]$ difiere de la ideal en un $0{,}006$\,\%. En las amplitudes mayores
(donde $\abs{A_V}V_p>\Vsat$) el cociente deja de valer $10$ porque la salida se recorta (cuestión~2).
\end{respuesta}

\begin{pregunta}[Recorte o \emph{clipping}]
¿Qué ocurre en la forma de la salida si la amplitud de la entrada se aumenta por encima de un cierto nivel?
Medir los niveles de entrada y salida para los que el fenómeno empieza.
\end{pregunta}

\begin{respuesta}
La salida del AO no puede superar las tensiones de saturación, algo menores que la alimentación:
$-V_{\mathrm{sat}}^{-}\le\vout\le V_{\mathrm{sat}}^{+}$. La relación lineal $\vout=A_V\vin$ solo vale mientras $\abs{A_V}\,\abs{\vin}$
no alcanza esos valores; por encima, las crestas de la sinusoide de salida se recortan en tramos planos.
El recorte empieza cuando
\begin{equation}\label{eq:i-umbral}
  \abs{A_V}\,V_{p,\mathrm{umbral}}=\Vsat\quad\Longrightarrow\quad V_{p,\mathrm{umbral}}=\frac{\Vsat}{\abs{A_V}} .
\end{equation}
Medimos $V_{\mathrm{sat}}^{+}=$ \dato\ y $V_{\mathrm{sat}}^{-}=$ \dato, y el recorte comienza para una amplitud de entrada
$V_{p,\mathrm{umbral}}=$ \dato, frente a $\Vsat/\abs{A_V}=$ \dato\ según \eqref{eq:i-umbral}. Las dos saturaciones pueden
no ser iguales, porque las etapas de salida positiva y negativa del 741 no son simétricas.
\end{respuesta}

\begin{pregunta}[Característica de transferencia]
Con la amplitud de entrada máxima, trazar $\Vout(\Vin)$ en modo X--Y. Señalar para qué amplitud de entrada
(en la representación temporal) se produce el corte y relacionarlo con lo medido en la representación X--Y.
\end{pregunta}

\begin{respuesta}
En modo X--Y el osciloscopio dibuja $\vout$ frente a $\vin$ eliminando el tiempo, de modo que vemos
directamente la característica de transferencia. Esperamos tres tramos:
\begin{equation}\label{eq:i-caract}
  \vout=\begin{cases}
    +V_{\mathrm{sat}}^{+}, & \vin\le-V_{\mathrm{sat}}^{+}/\abs{A_V},\\[2pt]
    A_V\,\vin, & \text{en la zona lineal},\\[2pt]
    -V_{\mathrm{sat}}^{-}, & \vin\ge V_{\mathrm{sat}}^{-}/\abs{A_V},
  \end{cases}
\end{equation}
una recta de pendiente $A_V$ por el origen y dos tramos horizontales. La pendiente medida es \dato\ y los codos
están en $\vin=$ \dato\ y \dato.
\begin{center}
\hueco{0.6\textwidth}{Característica $\vout(\vin)$ en modo X--Y}
\end{center}
Las abscisas de los codos son exactamente las amplitudes de entrada a las que empieza el recorte en la
representación temporal \eqref{eq:i-umbral}: ambas medidas describen el mismo fenómeno, la entrada del AO en
saturación, vista en un caso frente al tiempo y en el otro frente a $\vin$.
\end{respuesta}

\begin{pregunta}[Respuesta a $100\,\mathrm{kHz}$]
¿Qué ocurre en la salida si la frecuencia de la entrada se aumenta a $100\,\mathrm{kHz}$? Representar las ondas de
entrada y salida con sus puntos más significativos. ¿Cuánto vale la ganancia?
\end{pregunta}

\begin{respuesta}
A esta frecuencia deja de valer la hipótesis $A\to\infty$, y aparecen dos limitaciones del AO real.

\textbf{Ancho de banda finito.} La ganancia en lazo abierto del 741 decae como $\abs{A}\simeq f_t/f$, con
$f_t\simeq1\,\mathrm{MHz}$. Con realimentación, el inversor se comporta como un paso bajo de primer orden de frecuencia
de corte
\begin{equation}\label{eq:i-f3dB}
  f_{3\mathrm{dB}}=\frac{f_t}{1+R_2/R_1}\simeq\frac{1\,\mathrm{MHz}}{11}\simeq91\,\mathrm{kHz},
\end{equation}
de modo que, a $f=100\,\mathrm{kHz}$,
\begin{equation}\label{eq:i-100k}
  \abs{A_V}=\frac{R_2/R_1}{\sqrt{1+(f/f_{3\mathrm{dB}})^{2}}}\simeq6{,}7,\qquad
  \Delta\varphi=-\arctg\frac{f}{f_{3\mathrm{dB}}}\simeq-48^{\circ} :
\end{equation}
la ganancia baja y la salida se retrasa, además de la inversión, unos $48^{\circ}$, es decir,
$\Delta t=\abs{\Delta\varphi}/(360^{\circ}f)\simeq1{,}3\,\mu\mathrm{s}$.

\textbf{Velocidad de subida.} La pendiente de la salida no puede superar $\mathrm{SR}\simeq0{,}5\,\mathrm{V}/\mu\mathrm{s}$. Una
sinusoide de amplitud $V_p$ y frecuencia $f$ tiene pendiente máxima $2\pi fV_p$, luego solo se reproduce sin
deformarse si
\begin{equation}\label{eq:i-slew}
  V_{p,\mathrm{out}}\le\frac{\mathrm{SR}}{2\pi f}\simeq0{,}8\,\mathrm{V}\quad(f=100\,\mathrm{kHz});
\end{equation}
por encima, la salida se convierte en una onda triangular de pendientes $\pm\mathrm{SR}$.

\textbf{Medidas.} Observamos \dato. La ganancia medida es $\abs{A_V}=$ \dato, con un retraso adicional
$\Delta t=$ \dato; los máximos de la salida valen \dato\ y aparecen en $t=$ \dato.
\begin{center}
\hueco{0.6\textwidth}{Oscilograma a $100\,\mathrm{kHz}$, con máximos, mínimos\\y sus posiciones temporales}
\end{center}
\end{respuesta}

\begin{pregunta}[Valor de las resistencias]
¿Sería razonable elegir resistencias del orden de decenas de $\Omega$? ¿Y del orden de $\mathrm{M}\Omega$?
\end{pregunta}

\begin{respuesta}
No. El modelo ideal solo fija el cociente $R_2/R_1$, pero en el laboratorio ni los valores muy pequeños ni
los muy grandes funcionan bien. Veamos qué ocurriría en la práctica en cada caso.

\textbf{Decenas de ohmios} (por ejemplo, $R_1=10\,\Omega$, $R_2=100\,\Omega$).
\begin{itemize}
  \item \emph{El AO no puede dar tanta corriente.} Para sacar $5\,\mathrm{V}$ sobre $R_2=100\,\Omega$ haría falta
        \begin{equation}\label{eq:r-Iout}
          i=\frac{\vout}{R_2}=\frac{5\,\mathrm{V}}{100\,\Omega}=50\,\mathrm{mA},
        \end{equation}
        y el 741 limita su salida a unos $25\,\mathrm{mA}$. En el osciloscopio veríamos la sinusoide de salida
        recortada mucho antes de llegar a $\pm\Vsat$, con una amplitud menor que la esperada.
  \item \emph{No se quema, pero se calienta.} El 741 tiene protección interna contra cortocircuito, así que no
        se destruye; pero mientras limita disipa del orden de $0{,}1$--$0{,}25\,\mathrm{W}$, se nota caliente y sus
        parámetros (por ejemplo, el \emph{offset}) derivan.
  \item \emph{Las resistencias sí pueden sufrir.} $R_2$ disiparía hasta $V^{2}/R=25/100=0{,}25\,\mathrm{W}$ de pico, justo
        el límite de una resistencia de $1/4\,\mathrm{W}$; con valores aún menores se quemaría.
  \item \emph{Se pierde señal en el generador.} La entrada del inversor es $R_1=10\,\Omega$, más pequeña que la
        resistencia interna del generador ($50\,\Omega$): el generador marcaría una amplitud, pero al circuito
        llegaría solo una sexta parte.
\end{itemize}

\textbf{Megaohmios} (por ejemplo, $R_1=1\,\mathrm{M}\Omega$, $R_2=10\,\mathrm{M}\Omega$).
\begin{itemize}
  \item \emph{La señal se confunde con las corrientes parásitas del AO.} Con $\vin=0{,}1\,\mathrm{V}$, la corriente de señal es
        \begin{equation}\label{eq:r-IB}
          i=\frac{\vin}{R_1}=\frac{0{,}1\,\mathrm{V}}{1\,\mathrm{M}\Omega}=100\,\mathrm{nA},
        \end{equation}
        del mismo orden que la corriente que realmente entra por las entradas del 741 ($I_B\simeq80\,\mathrm{nA}$).
        En la práctica veríamos la salida desplazada en continua: $I_BR_2\simeq0{,}8\,\mathrm{V}$ de \emph{offset}, y la
        ganancia medida ya no sería $R_2/R_1$.
  \item \emph{Más ruido e interferencias.} Los nudos de tanta resistencia actúan como antenas: la señal saldría
        con ruido y con el zumbido de $50\,\mathrm{Hz}$ de la red.
  \item \emph{Pérdida de ancho de banda.} Las capacidades parásitas del montaje (unos pocos pF) junto con
        $R_2=10\,\mathrm{M}\Omega$ forman un filtro paso bajo con corte en torno a $10\,\mathrm{kHz}$: la ganancia empezaría a
        caer cerca de la frecuencia de trabajo.
\end{itemize}

\textbf{Conclusión.} Los valores razonables están entre $1$ y $100\,\mathrm{k}\Omega$, como los de nuestro diseño
($1$ y $10\,\mathrm{k}\Omega$): corrientes de miliamperios o menos, que el AO entrega sin problema, y muy por encima
de sus corrientes de polarización.
\end{respuesta}

% =========================================================================
\section{Circuito 2: amplificador no inversor}\label{sec:no-inversor}
% =========================================================================
Montamos el no inversor de la \cref{fig:no-inversor}, con $1+R_2/R_1=11$, a $1\,\mathrm{kHz}$.

% Figura: amplificador no inversor (circuito 2) con tensiones y corrientes del análisis.
\begin{figure}[ht]
\centering
\begin{circuitikz}[circuito]
  \draw (0,0) node[op amp, ao, noinv input up] (A) {};
  % nudo inversor (abajo) y red de realimentación
  \draw (A.out) -- ++(0.7,0) coordinate (O);
  \draw (A.-) -- ++(-0.6,0) coordinate (N) -- ++(0,-1.0) coordinate (M);
  \draw (M) -- ++(-1.0,0) coordinate (G1) to[R, l_=$R_1$, i>^=$i_1$] ++(0,-1.7) node[ground] {};
  \draw (M) to[R, l_=$R_2$, i<_=$i_2$] (M -| O) -- (O);
  % generador en la entrada no inversora
  \draw (A.+) to[short, -o] ++(-2.0,0) coordinate (Ein) node[above, tens] {$\vin$};
  \draw (Ein) -- ++(-0.8,0) to[sV, l_=$v_g$] ++(0,-2.0) node[ground] {};
  % carga y salida
  \draw (O) to[short, *-o] ++(1.0,0) node[right, tens] {$\vout$};
  \draw (O) ++(0.5,0) to[R, l=$R_L$] ++(0,-2.0) node[ground] {};
  \fill (M) circle (1.6pt);
  \node[tens, anchor=east, font=\scriptsize] at ($(N)!0.5!(M)+(-0.08,0)$) {$\vm=\vin$};
\end{circuitikz}
\caption{Amplificador no inversor (figura~3 del guion), con $R_1=1\,\mathrm{k}\Omega$,
$R_L=10\,\mathrm{k}\Omega$ y la alimentación omitida. La entrada $+$ no absorbe corriente; el nudo
inversor sigue a la entrada ($\vm=\vin$) y $R_1$--$R_2$ forman un divisor de $\vout$.}
\label{fig:no-inversor}
\end{figure}


\begin{pregunta}[Formas de onda y ganancia]
Dibujar las formas de onda de $\Vin=V_g$ y $\Vout$ detallando las escalas, y comparar la relación de
amplitudes o tensiones pico a pico con la $A_V$ buscada.
\end{pregunta}

\begin{respuesta}
Ahora la ganancia teórica y la experimental son
\begin{equation}\label{eq:n-Av}
  A_V^{\mathrm{teo}}=1+\frac{R_2}{R_1},\qquad
  A_V^{\mathrm{exp}}=+\frac{V_{\mathrm{out},pp}}{V_{\mathrm{in},pp}} ,
\end{equation}
con signo positivo porque la salida está en fase con la entrada.
\begin{center}\small
\renewcommand{\arraystretch}{1.3}
\begin{tabular}{@{}cccccc@{}}
\toprule
$V_p$ & Escala X, Y (V/div) & Base de tiempos & $V_{\mathrm{in},pp}$ & $V_{\mathrm{out},pp}$ & $A_V^{\mathrm{exp}}$\\
\midrule
\dato & \dato & \dato & \dato & \dato & \dato\\
\dato & \dato & \dato & \dato & \dato & \dato\\
\bottomrule
\end{tabular}
\end{center}
\begin{center}
\hueco{0.6\textwidth}{Oscilograma $\vin$, $\vout$ (en fase)}
\end{center}
Obtenemos $A_V^{\mathrm{exp}}=$ \dato\ frente a $A_V^{\mathrm{teo}}=$ \dato, una desviación de \dato. A diferencia del inversor,
la entrada no absorbe corriente ($\Zin\to\infty$), de modo que $\vin$ coincide con la tensión del generador.
\end{respuesta}

\begin{pregunta}[Característica de transferencia y recorte]
Dibujar la característica $\Vout(V_g)$. Señalar para qué amplitud de entrada (en la representación temporal) se
produce el recorte y relacionarlo con lo medido en la representación X--Y.
\end{pregunta}

\begin{respuesta}
La característica es una recta que pasa por el origen con pendiente $+11$ (la ganancia), hasta que la salida
llega a la tensión de saturación $\Vsat$; a partir de ahí la salida ya no puede crecer y la curva se vuelve
horizontal.

La salida empieza a recortarse cuando $11\cdot V_p$ alcanza $\Vsat$, es decir, para una amplitud de entrada
\begin{equation}\label{eq:n-umbral}
  V_{p,\mathrm{umbral}}=\frac{\Vsat}{A_V}=\frac{\Vsat}{11} .
\end{equation}
Medimos $V_{p,\mathrm{umbral}}=$ \dato\ en la pantalla en función del tiempo, y en el modo X--Y los codos de la
característica aparecen en $\vin=$ \dato\ y \dato. Ambas medidas deben coincidir, porque son el mismo
fenómeno visto de dos formas: el punto en que la recta se dobla en el modo X--Y es justo la amplitud a la
que la sinusoide de salida empieza a aplanarse.
\begin{center}
\hueco{0.6\textwidth}{Característica $\vout(\vin)$ en modo X--Y}
\end{center}
\end{respuesta}

\begin{pregunta}[Definición de la ganancia]
A la vista de las características de transferencia de ambos circuitos, discutir si sería conveniente o
posible definir la ganancia como $A_V=\partial\Vout/\partial\Vin$.
\end{pregunta}

\begin{respuesta}
Las dos características son lineales a trozos: una recta por el origen en la zona lineal y tramos
horizontales en saturación. La definición como cociente, $A_V=\vout/\vin$, solo tiene sentido en la zona
lineal; en saturación el cociente varía con la amplitud de entrada y no describe el circuito. La
definición incremental,
\begin{equation}\label{eq:n-deriv}
  A_V=\frac{\partial\vout}{\partial\vin}=
  \begin{cases}
    -R_2/R_1 \ \text{(inversor)},\quad 1+R_2/R_1 \ \text{(no inversor)}, & \text{zona lineal},\\[2pt]
    0, & \text{saturación},
  \end{cases}
\end{equation}
es la pendiente de la característica en cada punto. En la zona lineal coincide con el cociente, porque la
recta pasa por el origen; pero es más general: describe la respuesta a pequeñas variaciones en torno a cualquier
punto de trabajo, sigue siendo válida si hubiera una tensión de \emph{offset} (la recta no pasaría por el origen
y el cociente variaría con $\vin$) y muestra explícitamente que en saturación el amplificador deja de
amplificar. Por eso sí es conveniente, y es la definición que se usa para el comportamiento en pequeña señal.
\end{respuesta}

% =========================================================================
\section{Circuito 3: integrador y derivador inversores}\label{sec:int-der}
% =========================================================================
Montamos el integrador de la \figpanel{fig:int-der}{a} con los sets~1 ($R_F=100\,\mathrm{k}\Omega$) y~2
($R_F\to\infty$), y el derivador de la \figpanel{fig:int-der}{b} con $R=1\,\mathrm{k}\Omega$ y $C=1{,}59\,\mathrm{nF}$. Empezamos
por la respuesta temporal (cuestión~9 del informe, última del guion), porque de ella se deducen las demás.

% Figura: integrador inversor (figura 5a del guion) y derivador inversor (figura 5b).
\begin{figure}[ht]
\centering
\begin{minipage}[b]{0.48\textwidth}\centering
\begin{circuitikz}[circuito, scale=0.85, transform shape]
  \draw (0,0) node[op amp, ao] (A) {};
  \draw (A.-) -- ++(-0.5,0) coordinate (N);
  \draw (N) to[R, l_=$R$, i<^=$i$] ++(-2.2,0) to[short, -o] ++(-0.3,0) node[left, tens] {$\vin$};
  \draw (A.out) -- ++(0.6,0) coordinate (O);
  \draw (N) -- ++(0,1.2) coordinate (T1) to[C, l_=$C$] (T1 -| O) -- (O);
  \draw (T1) -- ++(0,1.0) coordinate (T2) to[R, l=$R_F$] (T2 -| O) -- (T1 -| O);
  \fill (T1) circle (1.6pt); \fill (T1 -| O) circle (1.6pt);
  \draw (O) to[short, *-o] ++(0.6,0) node[right, tens] {$\vout$};
  \draw (A.+) -- ++(-0.5,0) -- ++(0,-0.4) node[ground] {};
  \fill (N) circle (1.6pt);
\end{circuitikz}\\[2pt]
{\small (a)}
\end{minipage}\hfill
\begin{minipage}[b]{0.48\textwidth}\centering
\begin{circuitikz}[circuito, scale=0.85, transform shape]
  \draw (0,0) node[op amp, ao] (A) {};
  \draw (A.-) -- ++(-0.5,0) coordinate (N);
  \draw (N) to[C, l_=$C$, i<^=$i$] ++(-2.2,0) to[short, -o] ++(-0.3,0) node[left, tens] {$\vin$};
  \draw (A.out) -- ++(0.6,0) coordinate (O);
  \draw (N) -- ++(0,1.2) coordinate (T1) to[R, l=$R$] (T1 -| O) -- (O);
  \draw (O) to[short, *-o] ++(0.6,0) node[right, tens] {$\vout$};
  \draw (A.+) -- ++(-0.5,0) -- ++(0,-0.4) node[ground] {};
  \fill (N) circle (1.6pt);
\end{circuitikz}\\[2pt]
{\small (b)}
\end{minipage}
\caption{(a) Integrador inversor (figura~5a del guion); con $R_F\to\infty$ (set~2) es un integrador ideal.
(b) Derivador inversor (figura~5b del guion).}
\label{fig:int-der}
\end{figure}


\begin{pregunta}[Respuesta temporal]
Analizar la respuesta temporal de los circuitos de las figuras~5a (con $R_F=\infty$) y~5b: obtener $V_0(t)$ en
función de los parámetros del circuito y de la entrada $V_i(t)$.
\end{pregunta}

\begin{respuesta}
En ambos circuitos la entrada $+$ está a tierra, de modo que por (R2) el nudo inversor es una tierra virtual,
$\vm=0$, y por (R1) toda la corriente $i$ de la rama de entrada recorre la de realimentación.

\textbf{Integrador} ($R_F\to\infty$). La corriente por $R$ es $i=\vin/R$. En el condensador, la tensión entre sus
bornes es $\vm-\vout=-\vout$, y su corriente, $i=C\,\dif(\vm-\vout)/\dif t$:
\begin{equation}\label{eq:int1}
  \frac{\vin(t)}{R}=-C\,\frac{\dif\vout}{\dif t} .
\end{equation}
Integrando entre $0$ y $t$:
\begin{equation}\label{eq:int2}
  \resultado{\vout(t)=-\frac{1}{RC}\int_{0}^{t}\vin(\tau)\,\dif\tau+\vout(0)}
\end{equation}

\textbf{Derivador.} Ahora la corriente es la del condensador, con $\vin$ en un borne y la tierra virtual en el otro,
$i=C\,\dif\vin/\dif t$, y recorre $R$ hasta la salida, $\vm-\vout=Ri$:
\begin{equation}\label{eq:der}
  \resultado{\vout(t)=-RC\,\frac{\dif\vin}{\dif t}}
\end{equation}

\textbf{Set~1} ($R_F$ finita). El integrador tiene entonces $H(s)=-(R_F/R)/(1+sR_FC)$: un paso bajo de ganancia
$R_F/R=100$ en continua y corte en $f_c=1/(2\pi R_FC)\simeq1\,\mathrm{kHz}$, que solo integra para $f\gg f_c$. La
ventaja de $R_F$ es que limita la ganancia en continua: sin ella, cualquier \emph{offset} se integra y lleva la
salida a saturación.
\end{respuesta}

\begin{pregunta}[Integrador: elección de la señal]
¿Qué tipo de onda debemos seleccionar en el generador para que el osciloscopio permita verificar que la señal
de entrada es integrada por el circuito: la triangular o la rectangular?
\end{pregunta}

\begin{respuesta}
La \emph{rectangular}. Por \eqref{eq:int2}, en cada semiperiodo la entrada es constante, $\pm V_p$, y su integral es una
rampa: la salida es una onda triangular, cuya forma identifica sin ambigüedad la integración. Con una
triangular, la salida estaría formada por arcos de parábola, difíciles de distinguir a simple vista de una
sinusoide; y con una sinusoide solo obtendríamos otra sinusoide desfasada. Con el set~1, además, las rampas
solo son rectas si el semiperiodo es mucho menor que $R_FC\simeq159\,\mu\mathrm{s}$, es decir, si $f_0\gg f_c$.
\begin{center}
\hueco{0.47\textwidth}{Integrador, set 1:\\ entrada cuadrada y salida}\hfill
\hueco{0.47\textwidth}{Integrador, set 2:\\ entrada cuadrada y salida}
\end{center}
\end{respuesta}

\begin{pregunta}[Integrador: efecto de la frecuencia]
Manteniendo fijos el \emph{offset} y la amplitud de la entrada, variar su frecuencia $f_0$ entre $f_0/10$ y $10f_0$.
¿Influye la frecuencia en la ganancia en tensión del circuito?
\end{pregunta}

\begin{respuesta}
Sí. Para una entrada cuadrada de amplitud $V_p$ y periodo $T=1/f_0$, \eqref{eq:int2} da en cada semiperiodo una
rampa de excursión
\begin{equation}\label{eq:int-pp}
  V_{\mathrm{out},pp}=\frac{V_p}{RC}\,\frac{T}{2}=\frac{V_p}{2f_0RC},
  \qquad
  \frac{V_{\mathrm{out},pp}}{V_{\mathrm{in},pp}}=\frac{1}{4f_0RC},
\end{equation}
inversamente proporcional a la frecuencia: al multiplicar $f_0$ por $10$, la salida se divide por $10$.
\begin{center}\small
\renewcommand{\arraystretch}{1.3}
\begin{tabular}{@{}cccccc@{}}
\toprule
Set & $f_0$ & $V_{\mathrm{in},pp}$ & $V_{\mathrm{out},pp}$ & Cociente medido & $1/(4f_0RC)$\\
\midrule
2 & \dato & \dato & \dato & \dato & \dato\\
2 & \dato & \dato & \dato & \dato & \dato\\
2 & \dato & \dato & \dato & \dato & \dato\\
1 & \dato & \dato & \dato & \dato & \dato\\
\bottomrule
\end{tabular}
\end{center}
\end{respuesta}

\begin{pregunta}[Derivador: elección de la señal]
¿Qué tipo de onda debemos seleccionar para verificar que la señal de entrada es derivada por el circuito: la
triangular o la rectangular?
\end{pregunta}

\begin{respuesta}
La \emph{triangular}. Por \eqref{eq:der}, cada rampa de pendiente constante $\pm m$ da una salida constante
$\mp RCm$: la salida es una onda cuadrada. Con una rectangular, la derivada sería nula en los tramos planos e
infinita en los flancos: solo veríamos picos estrechos en cada transición, limitados por el AO real.
\begin{center}
\hueco{0.6\textwidth}{Derivador: entrada triangular y salida cuadrada}
\end{center}
\end{respuesta}

\begin{pregunta}[Derivador: efecto de la frecuencia]
Variando la frecuencia $f_0$ de la entrada entre $f_0/10$ y $10f_0$, ¿influye la frecuencia en la ganancia en
tensión del circuito?
\end{pregunta}

\begin{respuesta}
Sí, en sentido contrario al integrador. Una triangular de tensión pico a pico $V_{\mathrm{in},pp}$ y frecuencia $f_0$ sube
y baja $V_{\mathrm{in},pp}$ en cada semiperiodo, con pendiente $\abs{m}=2f_0V_{\mathrm{in},pp}$. Por \eqref{eq:der}, la salida es una
cuadrada de amplitud $RC\abs{m}$, es decir,
\begin{equation}\label{eq:der-pp}
  V_{\mathrm{out},pp}=2RC\abs{m}=4f_0RC\,V_{\mathrm{in},pp},
  \qquad
  \frac{V_{\mathrm{out},pp}}{V_{\mathrm{in},pp}}=4f_0RC ,
\end{equation}
proporcional a la frecuencia.
\begin{center}\small
\renewcommand{\arraystretch}{1.3}
\begin{tabular}{@{}ccccc@{}}
\toprule
$f_0$ & $V_{\mathrm{in},pp}$ & $V_{\mathrm{out},pp}$ & Cociente medido & $4f_0RC$\\
\midrule
\dato & \dato & \dato & \dato & \dato\\
\dato & \dato & \dato & \dato & \dato\\
\dato & \dato & \dato & \dato & \dato\\
\bottomrule
\end{tabular}
\end{center}
\end{respuesta}

\begin{pregunta}[Coherencia con la respuesta en frecuencia]
¿Son congruentes los resultados anteriores, en cuanto a la dependencia de la amplitud de salida con la
frecuencia para una entrada de amplitud constante? Justifíquese.
\end{pregunta}

\begin{respuesta}
Sí. En el dominio de la frecuencia, el integrador ideal y el derivador tienen
\begin{equation}\label{eq:H-int-der}
  H_{\mathrm{int}}(s)=-\frac{1}{sRC},\quad \abs{H_{\mathrm{int}}(\jj\omega)}=\frac{1}{2\pi fRC};\qquad
  H_{\mathrm{der}}(s)=-sRC,\quad \abs{H_{\mathrm{der}}(\jj\omega)}=2\pi fRC ,
\end{equation}
consecuencia directa de que en Laplace integrar es dividir por $s$ y derivar, multiplicar por $s$. Los cocientes
temporales \eqref{eq:int-pp} y \eqref{eq:der-pp} tienen exactamente la misma dependencia, $1/f$ y $f$: solo cambia
la constante ($4$ en lugar de $2\pi$), porque una cuadrada o una triangular no son sinusoides puras, sino sumas de
armónicos que el circuito trata cada uno según \eqref{eq:H-int-der}. Ambas descripciones coinciden además en la
frecuencia a la que la ganancia es unidad, del orden de $1/(2\pi RC)\simeq100\,\mathrm{kHz}$ con los valores nominales.
\end{respuesta}


\end{document}
